\section{Awareness: The Decision Filter}
\label{sec:awareness}

%%%%%%%%%%%%%%%%%%%%%%%%%%%%%%%%%%%%%%%%%%%%%%%%%%%%%%%%%%%%%%%%%%%%%%%%%%%%%%%
% CHAPTER 11: The Awareness Function (AWA)
% What is the person THINKING about in the decision moment?
%%%%%%%%%%%%%%%%%%%%%%%%%%%%%%%%%%%%%%%%%%%%%%%%%%%%%%%%%%%%%%%%%%%%%%%%%%%%%%%

Chapter~\ref{sec:fepsde-utility} established what \emph{could} generate utility:
144 components across three categories (INU, KNU, IDN), six dimensions (FEPSDE),
four time horizons, and two valences. But this full welfare structure exists
only in theory. In any actual decision moment, people do not---cannot---consider
all 144 components simultaneously.

This chapter introduces the \textbf{Awareness Function}: the filter that determines
which utility components are salient at the moment of decision.

%%%%%%%%%%%%%%%%%%%%%%%%%%%%%%%%%%%%%%%%%%%%%%%%%%%%%%%%%%%%%%%%%%%%%%%%%%%%%%%
\subsection{The Problem: From Potential to Actual}
\label{subsec:awareness-problem}
%%%%%%%%%%%%%%%%%%%%%%%%%%%%%%%%%%%%%%%%%%%%%%%%%%%%%%%%%%%%%%%%%%%%%%%%%%%%%%%

\subsubsection{The Gap Between Q and Decision}

\paragraph{The Full Welfare Function.}
From Chapter~\ref{sec:fepsde-utility}, total welfare is:

\begin{equation}
Q = \sum_{i} \omega_i \cdot U^i + \sum_{i<j} \gamma_{ij} \cdot U^i \cdot U^j
\end{equation}

This represents the \emph{comprehensive} evaluation---what would matter if
the person could consider everything.

\paragraph{The Decision Reality.}
But decisions happen in real time with limited attention:

\begin{itemize}
\item \textbf{Cognitive load:} Working memory holds 4-7 items
\item \textbf{Time pressure:} Many decisions are made in seconds
\item \textbf{Context cues:} Environment activates specific concerns
\item \textbf{Emotional state:} Current mood shapes what's salient
\end{itemize}

\paragraph{The Missing Link.}
Between the welfare function $Q$ and actual behavior lies a filter:

\begin{equation}
Q_{full} \xrightarrow{\text{Awareness Filter}} Q_{effective} \xrightarrow{\text{Decision}} \text{Behavior}
\end{equation}

This filter is the Awareness Function.

\subsubsection{Illustrative Examples}

\paragraph{Example 1: The Impulse Purchase.}
A person sees an attractive item on sale:

\begin{table}[htbp]
\centering
\caption{Impulse Purchase: Full vs. Activated Components}
\label{tab:impulse-purchase}
\renewcommand{\arraystretch}{1.3}
\begin{tabular}{@{}lcc@{}}
\toprule
\textbf{Component} & \textbf{In Full $Q$} & \textbf{Activated?} \\
\midrule
$U^{INU}_{E,0}$ (pleasure now) & Yes & \textbf{YES} \\
$U^{INU}_{F,0}$ (cost now) & Yes & Partially \\
$U^{INU}_{F,2}$ (budget impact) & Yes & No \\
$U^{KNU}_{F}$ (family budget) & Yes & No \\
$U^{IDN}$ (``I'm frugal'') & Yes & No \\
\bottomrule
\end{tabular}
\end{table}

The \emph{decision} reflects only what's activated, not the full $Q$.

\paragraph{Example 2: The Diet Choice at Lunch.}
A person choosing between salad and burger:

\begin{table}[htbp]
\centering
\caption{Diet Choice: Activation by Context}
\label{tab:diet-choice}
\renewcommand{\arraystretch}{1.3}
\begin{tabular}{@{}lccc@{}}
\toprule
\textbf{Component} & \textbf{Alone} & \textbf{With Health-Conscious Friend} \\
\midrule
$U^{INU}_{E,0}$ (taste) & HIGH & Medium \\
$U^{INU}_{P,2}$ (health) & Low & \textbf{HIGH} \\
$U^{KNU}_{S}$ (social fit) & Low & \textbf{HIGH} \\
$U^{IDN}$ (``healthy person'') & Low & \textbf{HIGH} \\
\bottomrule
\end{tabular}
\end{table}

The \emph{same person} with the \emph{same preferences} makes different choices
because different components are activated.

\paragraph{Example 3: The Retirement Saving Decision.}
An employee choosing contribution rate:

\begin{table}[htbp]
\centering
\caption{Retirement Saving: Activation by Framing}
\label{tab:retirement-framing}
\renewcommand{\arraystretch}{1.3}
\begin{tabular}{@{}lcc@{}}
\toprule
\textbf{Component} & \textbf{``Reduce Take-Home''} & \textbf{``Secure Your Future''} \\
\midrule
$U^{INU}_{F,0}$ (less money now) & \textbf{HIGH} & Medium \\
$U^{INU}_{F,3}$ (retirement security) & Low & \textbf{HIGH} \\
$U^{INU}_{E,3}$ (peace of mind) & Low & \textbf{HIGH} \\
$U^{KNU}_{F}$ (family security) & Low & HIGH \\
\bottomrule
\end{tabular}
\end{table}

Framing changes activation, which changes decision---without changing
underlying preferences ($\omega$) or welfare components ($U$).

%%%%%%%%%%%%%%%%%%%%%%%%%%%%%%%%%%%%%%%%%%%%%%%%%%%%%%%%%%%%%%%%%%%%%%%%%%%%%%%
\subsection{The Awareness Function: Definition}
\label{subsec:awareness-definition}
%%%%%%%%%%%%%%%%%%%%%%%%%%%%%%%%%%%%%%%%%%%%%%%%%%%%%%%%%%%%%%%%%%%%%%%%%%%%%%%

\subsubsection{Formal Definition}

\begin{definition}[The Awareness Function]
\label{def:awareness-function}
The Awareness Function $\mathbf{A}$ is a vector of activation levels for each
utility category:

\begin{equation}
\mathbf{A} = (a_{INU}, a_{KNU}, a_{IDN}) \quad \text{where } a_i \in [0, 1]
\end{equation}

where:
\begin{itemize}
\item $a_i = 1$: Category $i$ is fully salient (person is thinking about it)
\item $a_i = 0$: Category $i$ is not salient (not in awareness)
\item $a_i \in (0,1)$: Partial activation
\end{itemize}
\end{definition}

\paragraph{Effective Welfare.}
The welfare that actually guides decision-making is:

\begin{equation}
\boxed{Q^{eff} = \sum_{i \in \{INU, KNU, IDN\}} a_i \cdot \omega_i \cdot U^i}
\label{eq:effective-welfare}
\end{equation}

Note: This is the \emph{simplified} form. The full form including complementarities:

\begin{equation}
Q^{eff} = \sum_{i} a_i \cdot \omega_i \cdot U^i + \sum_{i<j} a_i \cdot a_j \cdot \gamma_{ij} \cdot U^i \cdot U^j
\end{equation}

Complementarities only operate when \emph{both} categories are activated.

\subsubsection{The Crucial Distinction: $\omega$ vs. $a$}

\paragraph{Two Different Parameters.}
The framework now has two parameters that could be confused:

\begin{table}[htbp]
\centering
\caption{Weight ($\omega$) vs. Awareness ($a$)}
\label{tab:omega-vs-a}
\renewcommand{\arraystretch}{1.4}
\begin{tabular}{@{}lll@{}}
\toprule
\textbf{Property} & \textbf{Weight ($\omega_i$)} & \textbf{Awareness ($a_i$)} \\
\midrule
\textbf{Meaning} & How important IS this to me? & Am I THINKING about it now? \\
\textbf{Timescale} & Stable (trait) & Situational (state) \\
\textbf{Source} & Personality, values, culture & Context, framing, cues \\
\textbf{Changeability} & Slow (months/years) & Fast (seconds/minutes) \\
\textbf{Consciousness} & Often implicit & Can be implicit or explicit \\
\textbf{Intervention} & Education, value change & Nudges, framing, reminders \\
\bottomrule
\end{tabular}
\end{table}

\paragraph{Why Both Are Needed.}
Consider a devoted parent ($\omega_{KNU}$ high for family):

\begin{itemize}
\item In a family context: $a_{KNU}$ high → Family considerations guide behavior
\item In a rushed work context: $a_{KNU}$ low → May make decisions ignoring family impact
\item The \emph{value} ($\omega_{KNU}$) hasn't changed; the \emph{awareness} ($a_{KNU}$) has
\end{itemize}

\paragraph{The Product $a_i \cdot \omega_i$.}
Behavior is driven by the \emph{product}:

\begin{equation}
\text{Effective weight} = a_i \cdot \omega_i
\end{equation}

\begin{itemize}
\item High $\omega$, low $a$: Values exist but not activated → No behavioral impact
\item Low $\omega$, high $a$: Activated but not valued → Weak behavioral impact
\item High $\omega$, high $a$: Values activated → Strong behavioral impact
\end{itemize}

\subsubsection{Dimension-Level Awareness (Extended Form)}

\paragraph{Category vs. Dimension Awareness.}
The basic model uses category-level awareness ($a_{INU}$, $a_{KNU}$, $a_{IDN}$).
For more precision, awareness can be defined at the dimension level:

\begin{equation}
\mathbf{A}_{extended} = \{a^i_d\} \quad \text{for } i \in \{INU, KNU, IDN\}, d \in FEPSDE
\end{equation}

This gives $3 \times 6 = 18$ awareness parameters.

\paragraph{When to Use Which.}
\begin{itemize}
\item \textbf{Category-level (3 parameters):} Sufficient for most analyses
\item \textbf{Dimension-level (18 parameters):} For fine-grained intervention design
\end{itemize}

Example: A health campaign might specifically target $a^{INU}_P$ (individual
physical awareness) without affecting $a^{INU}_F$ (individual financial awareness).

%%%%%%%%%%%%%%%%%%%%%%%%%%%%%%%%%%%%%%%%%%%%%%%%%%%%%%%%%%%%%%%%%%%%%%%%%%%%%%%
\subsection{Determinants of Awareness}
\label{subsec:awareness-determinants}
%%%%%%%%%%%%%%%%%%%%%%%%%%%%%%%%%%%%%%%%%%%%%%%%%%%%%%%%%%%%%%%%%%%%%%%%%%%%%%%

\subsubsection{What Shapes Awareness?}

Awareness is determined by a combination of:
\begin{enumerate}
\item Context and environment
\item Framing and communication
\item Cognitive state
\item Emotional state
\item Social situation
\item Time pressure
\end{enumerate}

\subsubsection{Context and Environment}

\paragraph{Physical Cues.}
The environment activates related utility components:

\begin{table}[htbp]
\centering
\caption{Environmental Cues and Awareness Activation}
\label{tab:environmental-cues}
\renewcommand{\arraystretch}{1.3}
\begin{tabular}{@{}lll@{}}
\toprule
\textbf{Environment} & \textbf{Activates} & \textbf{Mechanism} \\
\midrule
Gym & $a^{INU}_P \uparrow$ & Physical health becomes salient \\
Church & $a_{KNU} \uparrow$, $a_{IDN} \uparrow$ & Community and values activated \\
Shopping mall & $a^{INU}_E \uparrow$, $a^{INU}_F \downarrow$ & Pleasure up, cost awareness down \\
Hospital & $a^{INU}_P \uparrow$, $a^{KNU}_P \uparrow$ & Health for self and family \\
Nature & $a_{Eco} \uparrow$ & Environmental awareness \\
\bottomrule
\end{tabular}
\end{table}

\paragraph{Social Presence.}
Who is present affects what's salient:

\begin{table}[htbp]
\centering
\caption{Social Presence and Awareness}
\label{tab:social-presence}
\renewcommand{\arraystretch}{1.3}
\begin{tabular}{@{}lll@{}}
\toprule
\textbf{Present} & \textbf{Effect on Awareness} & \textbf{Mechanism} \\
\midrule
Alone & $a_{INU} \uparrow$, $a_{KNU} \downarrow$ & Self-focus default \\
Family & $a_{KNU} \uparrow$ & Family welfare salient \\
Colleagues & $a_{IDN} \uparrow$ (professional) & Professional identity activated \\
Strangers observing & $a_{IDN} \uparrow$ (public self) & Self-presentation concerns \\
Authority figure & $a_{KNU} \uparrow$ (norms) & Norm compliance salient \\
\bottomrule
\end{tabular}
\end{table}

\subsubsection{Framing and Communication}

\paragraph{``For You'' vs. ``For Us'' Framing.}

\begin{table}[htbp]
\centering
\caption{Framing Effects on Awareness}
\label{tab:framing-awareness}
\renewcommand{\arraystretch}{1.3}
\begin{tabular}{@{}lll@{}}
\toprule
\textbf{Frame} & \textbf{Awareness Effect} & \textbf{Example} \\
\midrule
``For you'' & $a_{INU} \uparrow$ & ``Save money on your energy bill'' \\
``For us/all'' & $a_{KNU} \uparrow$ & ``Together we can reduce emissions'' \\
``People like you'' & $a_{IDN} \uparrow$ & ``Responsible parents choose...'' \\
``Your future self'' & $a_{INU,t>0} \uparrow$ & ``Imagine yourself at 70...'' \\
Loss frame & $\lambda$ effect + $a_P \uparrow$ & ``Don't lose your chance...'' \\
\bottomrule
\end{tabular}
\end{table}

\paragraph{Temporal Framing.}

\begin{table}[htbp]
\centering
\caption{Temporal Framing and Time-Horizon Awareness}
\label{tab:temporal-framing}
\renewcommand{\arraystretch}{1.3}
\begin{tabular}{@{}lll@{}}
\toprule
\textbf{Frame} & \textbf{Effect} & \textbf{Example} \\
\midrule
``Right now'' & $a_{t=0} \uparrow$ & ``Instant savings'' \\
``In 30 years'' & $a_{t=3} \uparrow$ (but weak) & ``Your retirement...'' \\
``When you're 70'' & $a_{t=3} \uparrow$ (vivid) & Aging simulation apps \\
``Daily cost'' & $a_{t=0} \uparrow$ & ``Just \$2/day'' \\
``Lifetime cost'' & $a_{t=3} \uparrow$ & ``\$50,000 over your lifetime'' \\
\bottomrule
\end{tabular}
\end{table}

\subsubsection{Cognitive State}

\paragraph{Cognitive Load Effects.}

\begin{table}[htbp]
\centering
\caption{Cognitive Load and Awareness}
\label{tab:cognitive-load}
\renewcommand{\arraystretch}{1.3}
\begin{tabular}{@{}lll@{}}
\toprule
\textbf{Cognitive State} & \textbf{Awareness Pattern} & \textbf{Mechanism} \\
\midrule
Low load (relaxed) & All $a_i$ can be high & Capacity for reflection \\
High load (stressed) & $a_{INU} \uparrow$, others $\downarrow$ & Default to self-interest \\
Ego depletion & $a_{INU,t=0} \uparrow$ & Present-focus, low control \\
\bottomrule
\end{tabular}
\end{table}

\paragraph{The System 1/System 2 Connection.}
Kahneman's dual-process framework maps to awareness:

\begin{table}[htbp]
\centering
\caption{Dual Process and Awareness}
\label{tab:dual-process}
\renewcommand{\arraystretch}{1.3}
\begin{tabular}{@{}lll@{}}
\toprule
\textbf{System} & \textbf{Typical Awareness Pattern} & \textbf{Behavior} \\
\midrule
System 1 (fast, automatic) & $a_{INU} \uparrow$, $a_{t=0} \uparrow$ & Impulsive, self-interested \\
System 2 (slow, deliberate) & All $a_i$ moderate & Reflective, values-consistent \\
\bottomrule
\end{tabular}
\end{table}

\subsubsection{Emotional State}

\paragraph{Emotions as Awareness Modulators.}

\begin{table}[htbp]
\centering
\caption{Emotional State and Awareness}
\label{tab:emotion-awareness}
\renewcommand{\arraystretch}{1.3}
\begin{tabular}{@{}lll@{}}
\toprule
\textbf{Emotion} & \textbf{Awareness Effect} & \textbf{Behavioral Tendency} \\
\midrule
Fear & $a_P \uparrow$ (threat-relevant) & Risk avoidance, protection \\
Anger & $a_{INU} \uparrow$, $a_{KNU} \downarrow$ & Self-assertion, aggression \\
Guilt & $a_{KNU} \uparrow$, $a_{IDN} \uparrow$ & Prosocial, reparative \\
Pride & $a_{IDN} \uparrow$ & Identity-consistent behavior \\
Love & $a_{KNU} \uparrow$ (specific others) & Other-focused, caring \\
Anxiety & $a_{t=0} \uparrow$, $a_{P} \uparrow$ & Present-focused, threat-sensitive \\
\bottomrule
\end{tabular}
\end{table}

\subsubsection{Time Pressure}

\paragraph{The Rush Effect.}
Time pressure systematically shifts awareness:

\begin{equation}
\text{Time pressure} \uparrow \quad \Rightarrow \quad a_{INU} \uparrow, \; a_{t=0} \uparrow, \; a_{KNU} \downarrow, \; a_{IDN} \downarrow
\end{equation}

Under time pressure:
\begin{itemize}
\item System 1 dominates
\item Immediate self-interest becomes primary
\item Collective and identity considerations fade
\item Long-term consequences are ignored
\end{itemize}

\subsubsection{Summary: Awareness Determinant Matrix}

\begin{table}[htbp]
\centering
\caption{Summary: Factors Affecting Awareness}
\label{tab:awareness-summary}
\small
\renewcommand{\arraystretch}{1.2}
\begin{tabular}{@{}lcccc@{}}
\toprule
\textbf{Factor} & $a_{INU}$ & $a_{KNU}$ & $a_{IDN}$ & $a_{t>0}$ \\
\midrule
Time pressure & $\uparrow\uparrow$ & $\downarrow$ & $\downarrow$ & $\downarrow\downarrow$ \\
Social observation & $\sim$ & $\uparrow$ & $\uparrow\uparrow$ & $\sim$ \\
``For you'' framing & $\uparrow\uparrow$ & $\downarrow$ & $\sim$ & $\sim$ \\
``For us'' framing & $\downarrow$ & $\uparrow\uparrow$ & $\sim$ & $\sim$ \\
Identity prime & $\sim$ & $\sim$ & $\uparrow\uparrow$ & $\sim$ \\
Cognitive load & $\uparrow$ & $\downarrow$ & $\downarrow$ & $\downarrow$ \\
Deliberation time & $\sim$ & $\uparrow$ & $\uparrow$ & $\uparrow$ \\
Future visualization & $\sim$ & $\sim$ & $\sim$ & $\uparrow\uparrow$ \\
Group context & $\downarrow$ & $\uparrow\uparrow$ & $\uparrow$ & $\sim$ \\
\bottomrule
\end{tabular}
\begin{flushleft}
\small\textit{Note:} $\uparrow\uparrow$ = strong increase, $\uparrow$ = moderate increase, 
$\sim$ = neutral, $\downarrow$ = decrease, $\downarrow\downarrow$ = strong decrease
\end{flushleft}
\end{table}

%%%%%%%%%%%%%%%%%%%%%%%%%%%%%%%%%%%%%%%%%%%%%%%%%%%%%%%%%%%%%%%%%%%%%%%%%%%%%%%
\subsection{Nudges as Awareness Manipulation}
\label{subsec:nudges-awareness}
%%%%%%%%%%%%%%%%%%%%%%%%%%%%%%%%%%%%%%%%%%%%%%%%%%%%%%%%%%%%%%%%%%%%%%%%%%%%%%%

\subsubsection{The Nudge Mechanism Revealed}

\paragraph{What Nudges Actually Do.}
From the EBF perspective, most nudges work by changing awareness ($a$),
not by changing preferences ($\omega$) or outcomes ($U$):

\begin{equation}
\text{Nudge} \Rightarrow \Delta a_i \Rightarrow \Delta Q^{eff} \Rightarrow \Delta \text{Behavior}
\end{equation}

\paragraph{Classic Nudges Reinterpreted.}

\begin{table}[htbp]
\centering
\caption{Classic Nudges as Awareness Interventions}
\label{tab:nudges-awareness}
\renewcommand{\arraystretch}{1.3}
\begin{tabular}{@{}lll@{}}
\toprule
\textbf{Nudge} & \textbf{Awareness Mechanism} & \textbf{Effect} \\
\midrule
Default enrollment & Status quo bias + $a_{t=0} \downarrow$ & Reduces need to think about now \\
Social proof & $a_{KNU} \uparrow$ & ``Others do this'' activates collective \\
Commitment device & $a_{IDN} \uparrow$ (future) & ``Person who committed'' identity \\
Salience increase & $a_{\text{specific}} \uparrow$ & Makes particular component visible \\
Simplification & Reduces cognitive load & More components can be activated \\
Feedback & $a_{\text{relevant}} \uparrow$ & Makes consequences salient \\
Pre-commitment & $a_{t>0} \uparrow$ (at commitment time) & Future self activated when calm \\
\bottomrule
\end{tabular}
\end{table}

\subsubsection{Why Nudges Are Often Fragile}

\paragraph{The Awareness Decay Problem.}
Nudges change $a$, which is situational:

\begin{equation}
\text{Nudge at } t_0 \Rightarrow a_i(t_0) \uparrow \quad \text{but} \quad a_i(t_1) \rightarrow \text{baseline}
\end{equation}

Without sustained intervention, awareness returns to baseline, and behavior reverts.

\paragraph{The $\omega$-$a$ Distinction for Nudge Design.}

\begin{table}[htbp]
\centering
\caption{Nudges vs. Education: $a$ vs. $\omega$}
\label{tab:nudge-education}
\renewcommand{\arraystretch}{1.3}
\begin{tabular}{@{}lll@{}}
\toprule
\textbf{Intervention} & \textbf{Changes} & \textbf{Duration} \\
\midrule
Nudge & $a$ (awareness) & Temporary (needs repetition) \\
Education & $\omega$ (preferences) & Lasting (but slow) \\
Identity work & $\omega_{IDN}$, $U^{IDN}$ & Lasting if successful \\
\bottomrule
\end{tabular}
\end{table}

\paragraph{Implication for Intervention Design.}
\begin{itemize}
\item \textbf{For quick effect:} Nudge ($\Delta a$)
\item \textbf{For lasting effect:} Change $\omega$ or create identity shift
\item \textbf{For maximum effect:} Combine both---nudge to start, then build habits/values
\end{itemize}

%%%%%%%%%%%%%%%%%%%%%%%%%%%%%%%%%%%%%%%%%%%%%%%%%%%%%%%%%%%%%%%%%%%%%%%%%%%%%%%
\subsection{Awareness in the EBF Framework}
\label{subsec:awareness-bcm}
%%%%%%%%%%%%%%%%%%%%%%%%%%%%%%%%%%%%%%%%%%%%%%%%%%%%%%%%%%%%%%%%%%%%%%%%%%%%%%%

\subsubsection{The Updated Model Flow}

The complete EBF flow with Awareness:

\begin{equation}
\Psi \rightarrow U \rightarrow \mathbf{A} \rightarrow Q^{eff} \rightarrow \text{Decision}
\end{equation}

\begin{enumerate}
\item \textbf{Context ($\Psi$):} Shapes parameters and baseline awareness
\item \textbf{Utility ($U$):} 144 components exist in potential
\item \textbf{Awareness ($\mathbf{A}$):} Filters which components are salient
\item \textbf{Effective welfare ($Q^{eff}$):} What actually guides decision
\item \textbf{Decision:} Behavior follows $Q^{eff}$, not full $Q$
\end{enumerate}

\subsubsection{Connecting to Willingness (Chapter 12 Preview)}

\paragraph{The WTC Formula.}
Willingness to Change (Chapter~\ref{sec:willingness}) builds on awareness:

\begin{equation}
WTC = \sum_{i} a_i \cdot \omega_i \cdot \Delta U^i
\end{equation}

\begin{itemize}
\item $\Delta U^i$: Utility difference from change (from Chapter 10)
\item $\omega_i$: How much person cares (stable preference)
\item $a_i$: Whether they're thinking about it (situational)
\end{itemize}

Change happens when WTC exceeds a threshold---and awareness determines 
which $\Delta U^i$ components enter the calculation.

\subsubsection{Connecting to Journey (Chapter 13 Preview)}

\paragraph{Journey Phases and Awareness.}
The behavioral change journey is partly an awareness journey:

\begin{table}[htbp]
\centering
\caption{Journey Phases and Awareness Levels}
\label{tab:journey-awareness}
\renewcommand{\arraystretch}{1.3}
\begin{tabular}{@{}lll@{}}
\toprule
\textbf{Phase} & \textbf{Awareness Pattern} & \textbf{Characteristic} \\
\midrule
Pre-contemplation & All $a \approx 0$ for issue & ``Not thinking about it'' \\
Contemplation & $a$ rising & ``Maybe I should...'' \\
Preparation & $a$ high for relevant components & ``I'm planning to...'' \\
Action & $a$ very high & ``I'm doing it'' \\
Maintenance & $a$ moderate, habituated & ``It's normal now'' \\
\bottomrule
\end{tabular}
\end{table}

\subsubsection{Connecting to Segments (Chapter 14 Preview)}

\paragraph{Segments Have Different Awareness Defaults.}
From Chapter~\ref{sec:segments}, segments differ in $\omega$-profiles.
They also differ in \emph{baseline awareness}:

\begin{table}[htbp]
\centering
\caption{Segment Awareness Defaults}
\label{tab:segment-awareness}
\renewcommand{\arraystretch}{1.3}
\begin{tabular}{@{}lccc@{}}
\toprule
\textbf{Segment} & \textbf{Default $a_{INU}$} & \textbf{Default $a_{KNU}$} & \textbf{Default $a_{IDN}$} \\
\midrule
Individualist & HIGH & Low & Low \\
Collectivist & Low & HIGH & Low \\
Identitarian & Low & Low & HIGH \\
Balanced & Medium & Medium & Medium \\
\bottomrule
\end{tabular}
\end{table}

Segments with high $\omega_{KNU}$ also tend to have high baseline $a_{KNU}$---
but the two can diverge, which creates intervention opportunities.

%%%%%%%%%%%%%%%%%%%%%%%%%%%%%%%%%%%%%%%%%%%%%%%%%%%%%%%%%%%%%%%%%%%%%%%%%%%%%%%
\subsection{Summary: The Awareness Function}
\label{subsec:awareness-summary}
%%%%%%%%%%%%%%%%%%%%%%%%%%%%%%%%%%%%%%%%%%%%%%%%%%%%%%%%%%%%%%%%%%%%%%%%%%%%%%%

\begin{tcolorbox}[colback=gray!5!white, colframe=gray!75!black, title=Chapter 11 Summary: Awareness]

\textbf{The Core Insight:}
Decisions are guided not by full welfare $Q$, but by effective welfare $Q^{eff}$,
filtered by awareness $\mathbf{A}$.

\textbf{The Awareness Function:}
$$\mathbf{A} = (a_{INU}, a_{KNU}, a_{IDN}) \quad \text{where } a_i \in [0,1]$$

\textbf{Effective Welfare:}
$$Q^{eff} = \sum_{i} a_i \cdot \omega_i \cdot U^i$$

\textbf{The $\omega$-$a$ Distinction:}
\begin{itemize}
\item $\omega_i$: How important IS this to me? (stable, trait)
\item $a_i$: Am I THINKING about it now? (situational, state)
\item Behavior depends on the product $a_i \cdot \omega_i$
\end{itemize}

\textbf{Determinants of Awareness:}
\begin{itemize}
\item Context and environment
\item Framing and communication
\item Cognitive load
\item Emotional state
\item Social presence
\item Time pressure
\end{itemize}

\textbf{Nudges Work Through Awareness:}
Most nudges change $a$, not $\omega$ or $U$. This is why they're fast but often fragile.

\textbf{Key Patterns:}
\begin{itemize}
\item Time pressure → $a_{INU} \uparrow$, $a_{KNU/IDN} \downarrow$
\item Social observation → $a_{KNU} \uparrow$, $a_{IDN} \uparrow$
\item Cognitive load → Default to $a_{INU}$, $a_{t=0}$
\item Deliberation → All $a_i$ can increase
\end{itemize}

\textbf{Integration with EBF:}
$$\Psi \rightarrow U \rightarrow \mathbf{A} \rightarrow Q^{eff} \rightarrow WTC \rightarrow \text{Behavior}$$

\end{tcolorbox}

%%%%%%%%%%%%%%%%%%%%%%%%%%%%%%%%%%%%%%%%%%%%%%%%%%%%%%%%%%%%%%%%%%%%%%%%%%%%%%%
% END OF CHAPTER 11
%%%%%%%%%%%%%%%%%%%%%%%%%%%%%%%%%%%%%%%%%%%%%%%%%%%%%%%%%%%%%%%%%%%%%%%%%%%%%%%
